\documentclass[english,twoside, openright, hidelinks, 11 pt, class=report,crop=false]{standalone}
\input{../../preamb}
\input{../bm}

\begin{document}
\newpage
\section{Addisjon \label{Addisjon}}	
\subsection*{Addisjon med mengder; å legge til}
Når vi har en mengde og skal legge til mer, bruker vi \outl{plusstegnet} 
%id{plussym}
\sym{$ + $}
%\id
\!\!. Har vi 2 og skal legge til 3, skriver vi 
%id{2plus3}
\[ 2+3=5 \]
%\id
\fig{plusm1}
Rekkefølgen vi legger sammen tallene på har ikke  noe å si; å starte med 2 og så legge til 3 er det samme som å starte med 3 og så legge til 2:
%id{3plus2}
\[ 3+2=5 \]
%\id
\fig{plusm1a}
\spr{
	Et addisjonsstykke består av to eller flere \outl{ledd}\index{ledd} og én \outl{sum}\index{sum}. I regnestykket
	%id{2plus3}
	\[ 2+3=5 \]
	%\id
	er både 2 og 3 ledd, mens 5 er summen.\vsk
	
	Vanlige måter å si 
	%id{2plus3inline}
	$ 2+3 $ 
	%\id
	på er
	\begin{itemize}
		\item ''2 pluss 3''
		\item ''2 addert med 3''
	 	\item ''2 og 3 lagt sammen''
	\end{itemize}
Det å legge sammen tall kalles også \outl{å summere}.
} \regv
\newpage
\reg[Addisjon er kommutativ \label{adkom}]{
Summen er den samme uansett rekkefølge på leddene.
}
\eks{ \vsb 
%id{pluscommutative}
\alg{
2+5 &= 7 =5+2  \vn
6+3 &=9=3+6
}
%\id
}

\subsection*{Addisjon på tallinja; vandring mot høyre}
På en tallinje vil addisjon med positive tall  innebære vandring \textsl{mot \\høyre}:\regv
\eks[1]{ \vs
%id{2plus7}
\[ 2+7=9 \]
%\id
\fig{plus2}
}
\eks[2]{ \vs
%id{4plus11}
	\[ 4+11=15 \]
%\id
\fig{plus3}
}
\info{Betydningen av 
%id{symeq}
\sym{$\bm=$}
%\id
}{
%id{plussym}
\sym{$ + $} 
%\id
gir oss muligheten til å uttrykke tall på mange forskjellige måter, for eksempel er 
%id{5eq2plus3}
$ {5=2+3} $
%\id 
og 
%id{5eq1plus4}
$ {5=1+4} $
%\id
\!. I denne sammenhengen vil 
%id{symeq}
\sym{$ = $}
%\id
bety ''har samme verdi som''. Dette gjelder også ved subtraksjon, multiplikasjon og divisjon, som vi skal se på i de neste tre seksjonene.

}

\section{Subtraksjon \label{Subtraksjon}}
\subsection*{Subtraksjon med mengder: Å trekke ifra}
Når vi har en mengde og tar bort en del av den, bruker vi \outl{minustegnet} \sym{$ - $}. Har vi 5 og skal ta bort 3, skriver vi
%id{5min3}
\[ 5-{\color{red} 3}=2 \]
%\id
\fig{min1c}

\spr{
	Et subtraksjonsstykke består av to eller flere \outl{ledd} \index{ledd} og én\\ \outl{differanse}\index{differanse}. I subtraksjonsstykket
	%id{5min3lang}
	\[  5-3=2 \] 
	%\id
	er både 5 og 3 ledd, mens 2 er differansen. \vsk
	
	Vanlige måter å si 
	%id{5min3inline}
	$ 5-3 $
	%\id
	 på er
	\begin{itemize}
		\item ''5 minus 3'' \\
		\item ''5 fratrekt 3''
		\item ''3 subtrahert fra 5''
	\end{itemize}
} \vsk \vsk

\info{En ny tolkning av 0}{
	Innledningsvis i denne boka nevnte vi at 0 kan tolkes som \\''ingenting''. Subtraksjon gir oss muligheten til å uttrykke 0 via andre tall. For eksempel er 
	%id{7min7}
	$ {7-7=0} $ 
	%\id
	og 
	%id{19min19}
	$ {19-19=0} $
	%\id
	\!\!.
} \vsk

\begin{comment}
\info{Subtraksjon er \textsl{ikkje} kommutativ}{
$ 5-3 $ er ikkje lik\footnote{kva $ 3-5 $ er skal vi sjå på i \hrs{Negtal}{kapittel}} $ 3-5 $. Hadde alle rekneartar vore kommutative ville ikkje reknerekkefølga
}
\end{comment}
\newpage
\subsection*{Subtraksjon på tallinja: Vandring mot venstre}
I \hrs{Addisjon}{seksjon} har vi sett at 
%id{plussym}
\sym{$ + $}
%\id 
(med positive tall) innebærer at vi skal gå \textsl{mot høyre} langs tallinja. Med 
%id{minsym}
\sym{$ - $}
%\id
gjør vi omvendt, vi går \textsl{mot \\venstre}\!\!
%id{arrow}
\footnote{I figurer med tallinjer vil rødfargede piler indikere at man starter ved pilspissen og vandrer til andre enden.}:
%\id
 \regv
\info{Merk}{
	I \textsl{Eksempel 1} og \textsl{Eksempel 2} under går vi i motsatt retning av den som pila peker i. Dette kan først virke litt rart, men spesielt i \refkap{Negtal} vil det lønne seg å tenke slik.
} 
\eks[1]{ \vs
	%id{6min4}
	\[ 6-4=2 \]
	%\id
	\fig{mint}
}
\eks[2]{ \vs
	%id{12min5}
	\[ 12-7=5 \]
	%\id
	\fig{mint2}
}
\info{Omvendte operasjoner}{
Figurene som viser addisjon og subtraksjon på tallinja, illustrerer hvorfor vi sier at disse to er \outl{omvendte operasjoner}\index{omvendt operasjon} av hverandre. 
Sagt med et eksempel: Siden $2+4=6$, så vet vi at $6-4=2$ og at $6-2=4$.
}
\info{Å kansellere ledd}{
Om vi har et uttrykk hvor tall med lik verdi er addert like mange ganger som det er subtrahert, sier vi at de \outl{kansellerer} hverandre. I uttrykket
%id{mincncla}
\[ 4+3-4+4+7-4+1 \]
%\id
har vi like mange 4-ere addert som vi har 4-ere subtrahert. Dermed kansellerer 4-erene  hverandre, og vi kan forenkle uttrykket vårt til
%id{mincnclb}
\[ 3+7+1 \]
%\id
\textsl{Obs!} I \refkap{Brok} vil ordet \textit{kansellering} ha en annen betydning. 
}

\section{Ganging \label{Gonging} }

\subsection*{Ganging med heltall; innledende definisjon}
Når vi legger sammen like tall, kan vi bruke \outl{gangetegnet} 
%id{timessym}
\sym{$ \cdot $}
%\id
for å skrive regnestykkene våre kortere: \regv
\eks[]{\vsb
	%id{firsttimes}
	\alg{
		4+4+4 &= 4\cdot 3 \vn
		8+8 &=8\cdot 2 \vn
		1+1+1+1+1&= 1\cdot5 
	}
	%\id
} \regv
\spr{
	Et gangestykke består av to eller flere \outl{faktorer}\index{faktor} og ett \outl{produkt}\index{produkt}. I gangestykket
	%id{4times12}
	\[ {4\cdot 3=12} \]
	%\id
	er 4 og 3 faktorer, mens 12 er produktet. \vsk
	
	Vanlige måter å si 
	%id{4times12inline}
	$ 4\cdot3 $
	%\id
	på er
	\begin{itemize}
		\item ''4 ganger 3''  \\
		\item ''4 ganget med 3''\\
		\item ''4 multiplisert med 3''
	\end{itemize}
	
	Mange nettsteder og bøker på engelsk bruker symbolet 
	%id{timesx}
	\sym{$ \times $}
	%\id
	i steden for 
	%id{timessym}
	\sym{$ \cdot $}
	\!. I de fleste programmeringsspråk er 
	%id{times*}
	\sym{*}
	%\id
	symbolet for multiplikasjon.
}

\newpage
%id{amounttimes}
\subsection*{Ganging av mengder} \label{gangmengd}
%\id
La oss nå bruke en figur for å se for oss gangestykket 
%id{2times3}
$ 2\cdot3 $
%\id
\!\!:
\fig{2t3}
Og så kan vi legge merke til produktet av 
%id{3times2}
$ 3\cdot 2 $
%\id
\!\!:
\fig{3t2}
\reg[Multiplikasjon er kommutativ \label{gangkom}]{
	Produktet er det samme uansett rekkefølge på faktorene.
}
\eks[]{\vsb
	%id{timescommute}
	\alg{
		3\cdot 4 &=12= 4\cdot 3 \vn
		6\cdot 7 &=42= 7\cdot6 \vn
		8\cdot 9 &=72= 9\cdot8
	}
	%\id
}

\subsection*{Ganging på tallinja}
Vi kan også bruke tallinja for å regne ut gangestykker. For eksempel kan vi finne hva 
%id{2times4}
$ 2\cdot4 $
%\id
 er ved å tenke slik:
%id{2times4lang}
\[\text{''} 2\cdot 4 \text{ betyr å vandre 2 plasser \textsl{mot høyre}, 4 ganger.}\text{''} \]
%\id
%id{2times4eq8}
\[ 2\cdot4=8 \]
%\id
\fig{2t4l}
Også tallinja kan vi bruke for å overbevise oss om at rekkefølgen av faktorene ikke har noe å si:
%id{4times2lang}
\[\text{''} 4\cdot 2 \text{ betyr å vandre 4 plasser \textsl{mot høyre}, 2 ganger.}\text{''} \]
%\id
%id{4times2eq8}
\[ 4\cdot2=8 \]
%\id
\fig{4t2l}

\subsection*{Endelig definisjon av ganging med positive heltall}
Det ligger kanskje nærmest å tolke ''2 ganger 3'' som ''3, 2 ganger''. Da er
%id{2times3lang}
\[ \text{''2 ganger 3''}=3+3 \] 
%\id
%id{amounttimesref}
På side \pageref{gangmengd}
%\id
presenterete vi 
%id{2times3}
$ 2\cdot3 $
%\id
\!\!, altså ''2 ganger 3'', som 
%id{2plus2plus2}
$ {2+2+2} $
%\id
\!\!. Med denne tolkningen vil 
%id{3plus3}
$ {3+3} $
%\id
svare til 
%id{3times2}
$ {3\cdot2} $
%\id
\!\!, men nettopp det at multiplikasjon er en kommutativ operasjon (\rref{gangkom}) gjør at den ene tolkningen ikke utelukker den andre; 
%id{2times3eq2plus2plus2}
$ {2\cdot3 =2+2+2} 
%\id
$ og 
%id{2times3eq3plus3}
$ {2\cdot3=3+3} $ 
%\id
er to uttrykk med samme verdi.\regv

\reg[Ganging som gjentatt addisjon \label{ganggjad2}]{
	Ganging med et positivt heltal kan uttrykkes som gjentatt addisjon.
}
\eks[1]{\vsb
	%id{repeatedaddition}
	\alg{
		4+4+4 &= 4\cdot 3=3+3+3+3 \vn
		8+8 &=8\cdot 2=2+2+2+2+2+2+2 \vn
		1+1+1+1+1&= 1\cdot5 =5
	}
	%\id
}
\info{Merk}{
	At ganging med positive heltal kan uttrykkes som gjentatt addisjon, utelukker ikke andre uttrykk. Det er ikke feil å skrive at 
	%id{2times3alt}
	$ {2\cdot 3=1+5} $
	%\id
	\!\!.
}

\section{Divisjon \label{Divisjon}}
%id{divsym}
\sym{:}
%\id
er \outl{divisjonstegnet}.
I praksis har divisjon tre forskjellige betydninger, her eksemplifisert ved regnestykket 
%id{12div3}
$ 12:3 $
%\id
\!\!:\regv

\reg[Divisjon sine tre betydninger]{ \vs
	%id{divinterpret}
	\begin{itemize}
		\item \textbf{Inndeling av mengder} \\
		$ 12:3 = \text{''Antallet i hver gruppe når 12 deles inn i 3 like}$\\
		\hspace{1.6cm}store grupper'' 
		\item \textbf{Antall ganger} \\
		$ 12:3=\text{''Antall ganger 3 går på 12''} $
		\item \textbf{Omvendt operasjon av multiplikasjon}\\
		$ 12:3=\text{''Tallet man må gange 3 med for å få 12''} $
	\end{itemize}
	%\id
} \regv

\spr{ \label{sprakdiv}
	Et divisjonsstykke består av en \outl{dividend}\index{dividend}, en \outl{divisor}\index{divisor} og en\\ \outl{kvotient}\index{kvotient}.	
	I divisjonstykket
	%id{12div3display}
	\[  {12:3=4} \]
	%\id
	er 12 dividenden, 3 er divisoren og 4 er kvotienten.\vsk
	
	Vanlige måter å uttale 
	%id{12div3}
	$ 12:3 $
	%\id
	på er
	\begin{itemize}
		\item ''12 delt med/på 3''
		\item ''12 dividert med/på 3''
		\item ''12 på 3''
	\end{itemize}
	I noen sammenhenger blir 
	%id{12div3}
	$ {12:3} $ 
	%\id
	kalt ''\outl{forholdet}\index{forhold} mellom 12 og 3''. Da er 4 \outl{forholdstallet}\index{forholdstal}. \vsk
	
	Ofte brukes 
	%id{slash}
	\sym{$ / $}
	%\id
	 i steden for 
	 %id{divsym}
	 \sym{$:$}
	 %\id
	 \!\!, spesielt i programmeringsspråk.
}
\newpage
\subsection*{Divisjon av mengder}
Regnestykket 
%id{12div3}
$ {12:3} $
%\id
forteller oss at vi skal dele 12 inn i 3 like store \\grupper:
\fig{del1}
Vi ser at hver gruppe inneholder 4 ruter, dette betyr at
%id{12div3display}
\[ 12:3=4 \]
%\id


\subsection*{Antall ganger}
\fig{del2}
3 går 4 ganger på 12, altså er 
%id{12div3eq4}
$ 12:3=4 $
%\id
\!\!.


\subsection*{Omvendt operasjon av multiplikasjon}
Vi har sett at hvis vi deler 12 inn i 3 like grupper, får vi 4 i hver gruppe. Altså er 
%id{12div3eq4}
$ 12:3=4$
%\id
\!\!. Om vi legger sammen igjen disse gruppene, får vi naturligvis 12: 
\fig{del1b}
Men dette er det samme som å gange 4 med 3. Altså;
om vi vet at 
%id{4times3eq12}
$ {4\cdot 3=12} $
%\id
\!\!, så vet vi også at 
%id{12div3eq4}
$ {12:3=4} $
\!\!. I tillegg vet vi da at 
%id{12div4eq3}
$ {12:4=3} $
%\id
\!\!. 
\fig{del1a}


\eks[1]{Siden 
	%id{6times3eq18}
	$ 6\cdot 3 = 18  $
	%\id
	\!\!, er\vs
	%id{6times3}
	\alg{
		18:6= 3\vn
		18:3=6
	}
	%\id
}
\eks[2]{
	Siden 
	%id{5times7eq35}
	$ 5\cdot 7 = 35  $
	%\id
	\!\!, er\vs
	%id{5times7}
	\alg{
		35:5= 7\vn
		35:7=5
	}
	%\id
}

\end{document}